\documentclass[11pt]{article}

\usepackage[margin=1in]{geometry}
\usepackage{amsmath}
\usepackage{amssymb}
\usepackage{enumitem}
\usepackage[hidelinks]{hyperref}
\usepackage{longtable}
\usepackage{booktabs}

\setlist[itemize]{leftmargin=*, itemsep=0.5em}
\setlist[enumerate]{leftmargin=*, itemsep=0.5em}

\newcommand{\file}[1]{\texttt{#1}}
\newcommand{\loc}[2]{\texttt{#1}, line(s) \texttt{#2}}

\title{SPH Implementation Audit: \texttt{Assets/Scripts/Simulation}}
\author{}
\date{}

\begin{document}

\maketitle
\tableofcontents

\section{Scope and Method}

This document audits the 3D fluid solver in
\texttt{Assets/Scripts/Simulation} against the theory collected in
\texttt{sph\_study\_notes.tex} (referred to below simply as the
\emph{study notes}). The reviewed source files are:

\begin{itemize}
    \item \file{FluidSim.cs} --- MonoBehaviour driving the simulation loop,
    buffer management, and per-frame parameter upload.
    \item \file{Spawner3D.cs} --- initial particle placement.
    \item \file{Compute/FluidSim.compute} --- all GPU kernels (external
    forces, spatial hashing, reordering, density, pressure, viscosity,
    position integration, density-texture rendering, and foam/spray/bubble
    particles).
    \item \file{Compute/FluidMaths3D.hlsl} --- smoothing-kernel
    definitions.
    \item \file{Compute/SpatialHash3D.hlsl} --- uniform-grid cell hashing
    helpers.
    \item \file{Assets/Scripts/Helpers/SpatialHash/SpatialHash.cs} ---
    (outside \texttt{Simulation}, but instantiated directly by
    \file{FluidSim.cs}) wraps the GPU counting sort used to build the
    neighbor lookup table.
\end{itemize}

For each SPH concept, the exact implementation is identified with file and
line references, and cross-referenced to the corresponding section of the
study notes. The final section lists mechanisms described in the study notes
(or standard in the wider SPH literature) that are \emph{not} present in this
codebase, together with suggested improvements.

\section{Spatial Discretization and Neighbor Search}

The study notes describe discretization abstractly (Section 4, particularly
\emph{What Is Discretization?}, study notes lines 540--550) but do not detail
a specific neighbor-search algorithm; they point to Ihmsen et al.\ (2014),
\emph{SPH Fluids in Computer Graphics}, as the reference survey for
neighborhood search (study notes lines 994--999).

The implementation uses a \textbf{uniform background grid with spatial
hashing}, the standard approach surveyed in that paper:

\begin{itemize}
    \item \textbf{Cell indexing.} \loc{SpatialHash3D.hlsl}{46--48}
    (\texttt{GetCell3D}) maps a world position to an integer cell coordinate
    by dividing by the smoothing radius, i.e.\ cell size equals the kernel
    support $h$, matching the usual convention that at most the 27 cells in a
    $3\times3\times3$ block around a particle can contain neighbors within
    $h$.
    \item \textbf{Hashing.} \loc{SpatialHash3D.hlsl}{53--63}
    (\texttt{HashCell3D}) hashes the cell coordinate using large prime
    multipliers (lines 41--43) combined with a block-wise wrap-around
    (block size 50) to reduce large-coordinate overflow issues, then
    \loc{SpatialHash3D.hlsl}{65--68} (\texttt{KeyFromHash}) reduces the hash
    modulo the particle count to produce a bounded table key.
    \item \textbf{27-cell neighbor stencil.} \loc{SpatialHash3D.hlsl}{1--39}
    defines the fixed \texttt{offsets3D[27]} array enumerating every
    neighboring cell offset in 3D.
    \item \textbf{Key generation kernel.} \loc{FluidSim.compute}{146--164}
    (\texttt{UpdateSpatialHash}) computes a key per particle from its
    \emph{predicted} position (see Section~\ref{sec:timeintegration}).
    \item \textbf{Sort and offset table.} The keys are sorted with a GPU
    counting sort and converted into a start-offset table by
    \texttt{SpatialHash.Run} (\file{Helpers/SpatialHash/SpatialHash.cs},
    lines 30--36), which calls \texttt{GPUCountSort} and
    \texttt{SpatialOffsetCalculator}. This produces, for each hashed key, the
    index at which that key's particle run begins in the sorted array ---
    the same ``counting sort into an offset table'' neighbor-search strategy
    discussed by Ihmsen et al.
    \item \textbf{Cache-coherent reordering.} \loc{FluidSim.compute}{165--185}
    (\texttt{Reorder} and \texttt{ReorderCopyBack}) physically rearrange
    \texttt{Positions}, \texttt{PredictedPositions}, and \texttt{Velocities}
    into sorted order so that later neighbor loops read contiguous memory
    instead of indirecting through \texttt{SortedIndices} on every access.
    This is a GPU-performance optimization with no direct analogue in the
    study notes, but it materially affects achievable particle counts.
    \item \textbf{Neighbor traversal.} The shared pattern used in
    \texttt{CalculateDensities}, \texttt{CalculatePressureForce},
    \texttt{CalculateViscosity}, \texttt{UpdateDensityTexture}, and
    \texttt{UpdateWhiteParticles} (e.g.\ \loc{FluidSim.compute}{187--226})
    loops over the 27 neighboring cells, looks up each cell's start offset,
    and walks forward through the sorted array until the key changes,
    discarding neighbors outside the squared smoothing radius.
\end{itemize}

\section{Smoothing Kernels}

The study notes derive and verify a single kernel in detail: the
\textbf{cubic B-spline ($M_4$) kernel} (Section 4.2.3, study notes lines
676--753), including its normalization constants $\sigma_1,\sigma_2,\sigma_3$
and its piecewise gradient (study notes lines 851--863). The implementation
does \emph{not} use the cubic spline. Instead, \file{FluidMaths3D.hlsl}
implements the kernel family popularized by M\"uller, Charypar, and Gross
(2003) (cited in study notes lines 950--955):

\begin{itemize}
    \item \textbf{Poly6}, \loc{FluidMaths3D.hlsl}{17--26}
    (\texttt{SmoothingKernelPoly6}): $W\propto(h^2-r^2)^3$. Used only by the
    viscosity kernel (\loc{FluidSim.compute}{343, 490--509}); see
    Section~\ref{sec:viscosity}.
    \item \textbf{Spiky, power 2}, \loc{FluidMaths3D.hlsl}{39--47}
    (\texttt{SpikyKernelPow2}) and its derivative
    \loc{FluidMaths3D.hlsl}{59--67} (\texttt{DerivativeSpikyPow2}): used as
    the \emph{density} kernel (\loc{FluidMaths3D.hlsl}{69--73},
    \texttt{DensityKernel}/\texttt{DensityDerivative}). Using a
    Spiky-type kernel (non-zero derivative as $r\to0$) for density, rather
    than the smoother Poly6 kernel that is more conventional for density
    summation, is a deliberate anti-clustering choice: Poly6's derivative
    vanishes at $r=0$, which lets particles approach arbitrarily closely
    without a growing density gradient to push them apart, whereas Spiky's
    non-vanishing derivative keeps the density gradient finite and repulsive
    at small separations.
    \item \textbf{Spiky, power 3}, \loc{FluidMaths3D.hlsl}{28--36}
    (\texttt{SpikyKernelPow3}) and its derivative
    \loc{FluidMaths3D.hlsl}{49--57} (\texttt{DerivativeSpikyPow3}): used as
    the \emph{near-density} kernel (\loc{FluidMaths3D.hlsl}{75--83},
    \texttt{NearDensityKernel}/\texttt{NearDensityDerivative}); see
    Section~\ref{sec:nearpressure}.
    \item \textbf{Linear kernel}, \loc{FluidMaths3D.hlsl}{8--15}
    (\texttt{LinearKernel}): defined but \emph{never called} anywhere in
    \file{FluidSim.compute}. Dead code / unused leftover.
\end{itemize}

All normalization constants (\texttt{K\_SpikyPow2}, \texttt{K\_SpikyPow3},
\texttt{K\_SpikyPow2Grad}, \texttt{K\_SpikyPow3Grad}) are declared
uninitialized in the shader (\loc{FluidMaths3D.hlsl}{3--6}) and are uploaded
once per frame from the CPU side by \texttt{UpdateSmoothingConstants}
(\loc{FluidSim.cs}{339--351}), which recomputes them only when
\texttt{smoothingRadius} changes. This mirrors the study notes' general
normalization procedure (study notes lines 716--730) but specialized
in closed form to the Spiky kernels rather than computed via the generic
integral $\int_0^1 f(q)q^{d-1}dq$.

Notably, the implementation never explicitly includes a particle mass $m_j$
in any kernel sum (see Section~\ref{sec:notimplemented}); all sums are of the
unweighted kernel value, which is only correct if every particle is assumed
to carry equal, implicitly-unit mass.

\section{Density and Near-Density Fields}

\loc{FluidSim.compute}{187--226} (\texttt{CalculateDensitiesAtPoint}) and
the wrapping kernel \loc{FluidSim.compute}{229--241}
(\texttt{CalculateDensities}) implement the \textbf{summation density}
approach described in the study notes,
\begin{equation*}
    \rho_i \approx \sum_j m_j W_{ij}
\end{equation*}
(study notes lines 822--827), but with $m_j\equiv1$ and $W_{ij}=$
\texttt{DensityKernel}. Alongside the ordinary density, a second
\textbf{near-density} field $\rho_i^{\text{near}}=\sum_j
\texttt{NearDensityKernel}(r_{ij},h)$ is accumulated in the same loop and
returned as the second component of the \texttt{float2} result
(\loc{FluidSim.compute}{222--224}). Both fields are evaluated using
\emph{predicted} positions, not the current positions (see
Section~\ref{sec:timeintegration}).

This dual density/near-density scheme is not discussed anywhere in the study
notes or its reading list; it is a lineage from Clavet, Lowe, and Beaudoin's
\emph{Particle-based Viscoelastic Fluid Simulation} (2005)
``double density relaxation'' technique, adapted here into a conventional
force-based (rather than position-relaxation-based) SPH solver. The study
notes' paper list (Section 5) jumps from M\"uller et al.\ 2003 directly to
Monaghan 2005 and Becker--Teschner 2007 (study notes lines 950--969) without
covering this intermediate near-pressure technique, so this is a method
present in the code but absent from the notes.

\section{Pressure Closure and Near-Pressure Stabilization}
\label{sec:nearpressure}

The study notes' Section 2.1.4 (\emph{Pressure Closure}, study notes lines
260--292) presents two closures: WCSPH with Tait's equation
$p=\frac{\rho_0c_0^2}{\gamma}\left[(\rho/\rho_0)^\gamma-1\right]$, and
incompressible/projection SPH with a pressure Poisson equation. The
implementation uses neither in full. Instead,
\loc{FluidSim.compute}{85--88} (\texttt{PressureFromDensity}) implements the
\emph{linearized} equation of state
\begin{equation*}
    p_i = (\rho_i - \rho_{\text{target}})\cdot k_{\text{pressure}},
\end{equation*}
i.e.\ Tait's equation with $\gamma=1$ and no reference speed of sound,
exposed as \texttt{targetDensity} and \texttt{pressureMultiplier}
(\loc{FluidSim.cs}{22--23}). This is the simplified stiffness relation used
in M\"uller et al.\ (2003) rather than the more general exponentiated Tait
form the study notes write out --- the notes describe the general WCSPH
family (study notes lines 437--443 also discuss the $\gamma$ exponent) but
the code implements only its simplest linear member.

A second, independent pressure-like term, \texttt{NearPressureFromDensity}
(\loc{FluidSim.compute}{90--93}), computes
$p_i^{\text{near}}=\rho_i^{\text{near}}\cdot k_{\text{near}}$
(\texttt{nearPressureMultiplier}, \loc{FluidSim.cs}{24}) from the
near-density field of Section 4. Both pressures are combined, symmetrized
between particle pairs, and converted into a force in
\loc{FluidSim.compute}{267--346} (\texttt{CalculatePressureForce}):
\begin{equation*}
    \mathbf{f}_i^{\text{pressure}} = -\sum_j
        \nabla_iW_{ij}\,\frac{p_i+p_j}{2\rho_j}
        -\sum_j
        \nabla_iW^{\text{near}}_{ij}\,\frac{p_i^{\text{near}}+p_j^{\text{near}}}{2\rho_j^{\text{near}}}
\end{equation*}
(signs folded into the kernel-derivative helper functions;
\loc{FluidSim.compute}{329--331}). Using the \emph{averaged} pressure
$(p_i+p_j)/2$ for the pair force is exactly the symmetric pairwise
formulation the study notes credit with exact momentum conservation (study
notes lines 104--107, \emph{Conservation and coupling}). The near-pressure
term is an additional short-range repulsive correction that prevents the
particle clustering/tensile-instability failure mode the study notes warn
about under ``Stability artifacts'' (study notes lines 126--131), without
needing the artificial-pressure or $\delta$-SPH corrections the notes list
as alternatives (study notes lines 459--463).

The resulting acceleration is applied directly to velocity
(\loc{FluidSim.compute}{347--349}), and an ad hoc aerodynamic drag term is
applied to particles with fewer than 8 neighbors
(\loc{FluidSim.compute}{351--355}) to damp free-flying/airborne particles
--- an engineering stabilization with no counterpart in the study notes.

\section{Viscosity}
\label{sec:viscosity}

The study notes' momentum equation includes a physical viscous term
$\mu\nabla^2\mathbf{v}$ (study notes lines 224--232), and Section 3.4
separately lists ``XSPH velocity smoothing'' as a distinct
\emph{regularization} technique, not a physical viscosity model (study notes
lines 464--468). The implementation's \texttt{CalculateViscosity} kernel
(\loc{FluidSim.compute}{409--459}) computes
\begin{equation*}
    \Delta\mathbf{v}_i = \texttt{viscosityStrength}\cdot\Delta t\sum_j
        (\mathbf{v}_j-\mathbf{v}_i)\,W^{\text{poly6}}(r_{ij},h)
\end{equation*}
and adds it directly to particle velocity (\loc{FluidSim.compute}{456--458}).
This is XSPH velocity averaging, not a discretized Laplacian of velocity: it
pulls each particle's velocity toward the kernel-weighted average of its
neighbors' velocities, rather than modeling a stress-divergence force scaled
by dynamic viscosity $\mu$. It therefore matches the regularization
technique named in study notes Section 3.4 rather than the physical-viscosity
discretizations referenced via Morris, Fox, and Zhu (1997) in the study
notes' paper list (study notes lines 940--944). The kernel is skipped
entirely when \texttt{viscosityStrength == 0} (\loc{FluidSim.cs}{332}), which
is also the default value (\loc{FluidSim.cs}{25}).

\section{Time Integration}
\label{sec:timeintegration}

The study notes' Section 3.5 lists explicit Euler, leapfrog, Verlet,
predictor--corrector, and Runge--Kutta as integrator choices with a CFL-type
timestep restriction (study notes lines 469--494). The implementation uses a
\textbf{semi-implicit (symplectic) Euler} scheme with position prediction and
fixed sub-stepping:

\begin{itemize}
    \item \textbf{Gravity and prediction}, \loc{FluidSim.compute}{133--144}
    (\texttt{ExternalForces}): velocity is updated by gravity
    ($\mathbf{v} \mathrel{+}= \mathbf{g}\,\Delta t$), and a
    \emph{predicted} position is computed as $\mathbf{x} + \mathbf{v}\cdot
    \frac1{120}$ (a fixed $1/120\,\text{s}$ look-ahead, independent of the
    actual timestep) rather than the true $\Delta t$. All density/pressure
    computations in this step then use these predicted positions, not the
    actual particle positions, so the pressure solve reacts to where
    particles are about to be rather than where they currently sit. This
    look-ahead/prediction trick is not discussed in the study notes.
    \item \textbf{Position integration}, \loc{FluidSim.compute}{461--475}
    (\texttt{UpdatePositions}): the real \texttt{Positions} buffer is
    advanced by the final velocity, $\mathbf{x}\mathrel{+}=\mathbf{v}\,\Delta
    t$, after which boundary collisions are resolved (Section
    \ref{sec:boundary}).
    \item \textbf{Fixed sub-stepping instead of a CFL formula.} Rather than
    computing $\Delta t$ from the stability formula in study notes lines
    476--486 (which involves $c_0$, $v_{\text{sig}}$, $a_{\max}$, and $\nu$),
    \texttt{RunSimulationFrame} (\loc{FluidSim.cs}{284--308}) simply divides
    the frame's delta time by a fixed, user-set
    \texttt{iterationsPerFrame} (\loc{FluidSim.cs}{18}) and runs that many
    identical sub-steps (\loc{FluidSim.cs}{290--293}). An optional frame-rate
    floor, \texttt{maxTimestepFPS} (\loc{FluidSim.cs}{17}), caps the
    per-frame $\Delta t$ so that a slow frame runs the simulation in slow
    motion instead of taking a single large, unstable step --- a crude but
    effective proxy for the notes' CFL safety coefficient
    $C_{\text{CFL}}$ (study notes lines 487--489), implemented without ever
    measuring particle velocity or acceleration.
\end{itemize}

\section{Boundary Handling}
\label{sec:boundary}

The study notes discuss boundary treatment at several points: as a
limitation of SPH in general (``Solid walls are not automatically
represented accurately,'' study notes lines 132--134), as a required initial/
boundary condition (Section 2.1.6, study notes lines 315--331), as a tunable
parameter category (Section 3.6, study notes lines 496--504), and via the
Adami--Hu--Adams (2012) generalized wall boundary condition in the paper list
(study notes lines 988--992). All of these describe \emph{boundary
particles} (ghost, dummy, or dynamic particles that participate in the SPH
sums).

The implementation does not use boundary particles at all. Collisions are
resolved purely geometrically in \texttt{ResolveCollisions}
(\loc{FluidSim.compute}{101--131}): the box is transformed into local
(unit-cube) space via \texttt{worldToLocal}, any axis whose local coordinate
exceeds $\pm0.5$ is clamped back to the boundary, and the corresponding
velocity component is negated and scaled by \texttt{collisionDamping}
(\loc{FluidSim.cs}{26}), then both are transformed back to world space. This
supports an arbitrarily translated/rotated/scaled rectangular container (via
\texttt{localToWorld}/\texttt{worldToLocal}, set from the \texttt{FluidSim}
transform in \loc{FluidSim.cs}{372--373}) at negligible cost, but it is a
purely kinematic, perfectly-reflecting-with-damping wall model: it has no
SPH pressure/viscosity interaction with wall particles, so no-slip,
free-slip, or partial-slip conditions (study notes lines 500--501) cannot be
modeled, and only convex rectangular-box containers are supported ---
arbitrary solid geometry (e.g.\ ramps, meshes, moving rigid bodies) is out of
scope for this routine.

Two uniform parameters, \texttt{edgeForce} and \texttt{edgeForceDst}
(\loc{FluidSim.compute}{44--45}), are declared but \emph{never read} anywhere
in \file{FluidSim.compute}. Likewise, the mouse/pointer interaction uniforms
\texttt{interactionInputPoint}, \texttt{interactionInputStrength}, and
\texttt{interactionInputRadius} (\loc{FluidSim.compute}{51--53}) are declared
but never referenced by any kernel. Both appear to be vestigial from a 2D
sibling implementation of this simulation (not present in the audited
\texttt{Simulation} folder) and represent dead/unused code paths here: there
is currently no user-interaction force (e.g., a mouse-driven push/pull) and
no additional soft edge-repulsion force beyond the hard
\texttt{ResolveCollisions} clamp.

\section{GPU-Specific Structure}

Beyond the numerical method itself, the solver is organized as 11 compute
kernels dispatched once per sub-step from \texttt{RunSimulationStep}
(\loc{FluidSim.cs}{323--337}):
\texttt{ExternalForces} $\to$ \texttt{UpdateSpatialHash} $\to$ (GPU sort)
$\to$ \texttt{Reorder} $\to$ \texttt{ReorderCopyBack} $\to$
\texttt{CalculateDensities} $\to$ \texttt{CalculatePressureForce} $\to$
(optionally) \texttt{CalculateViscosity} $\to$ \texttt{UpdatePositions}.
All buffers are plain \texttt{ComputeBuffer}s created and bound once in
\texttt{Initialize} (\loc{FluidSim.cs}{96--255}); there is no
dynamic/adaptive resizing of the particle count at runtime (the spawn count
is fixed by \file{Spawner3D.cs} at scene start, \loc{Spawner3D.cs}{20--33}).
This is an implementation-architecture detail rather than an SPH numerical
method, but it bounds what the solver can do (e.g., no particle
emitters/sinks, no adaptive refinement) independent of the physics used.

\section{Extensions Beyond Classical SPH: Foam, Spray, and Bubbles}

A substantial portion of \file{FluidSim.compute}
(\loc{FluidSim.compute}{356--607}, inside \texttt{CalculatePressureForce},
plus the \texttt{UpdateWhiteParticles} and
\texttt{WhiteParticlePrepareNextFrame} kernels,
\loc{FluidSim.compute}{490--607}) implements a secondary particle system for
visual foam/spray/bubble effects, following Ihmsen, Akinci, Gissler, and
Teschner's \emph{Unified Spray, Foam and Bubbles for Particle-Based Fluids}
(SIGGRAPH 2012) --- a paper not included in the study notes' reading list
(Section 5), even though it is a direct companion to the Ihmsen et al.\ 2014
survey that \emph{is} listed (study notes lines 994--999). Specifically:

\begin{itemize}
    \item \textbf{Trapped-air potential.} For each fluid particle, a
    ``weighted velocity difference'' is accumulated from converging
    neighbor velocity pairs (\loc{FluidSim.compute}{313--323}), then combined
    with a kinetic-energy factor to produce a probabilistic spawn count for
    new foam/spray particles (\loc{FluidSim.compute}{358--366}), sampled over
    a cylindrical volume swept along the particle's velocity
    (\loc{FluidSim.compute}{368--388}).
    \item \textbf{Classification.} Each foam/spray/bubble particle counts its
    fluid neighbors in \texttt{UpdateWhiteParticles}
    (\loc{FluidSim.compute}{509--540}) and is classified as spray (few
    neighbors), bubble (many neighbors), or foam (in between)
    (\loc{FluidSim.compute}{542--545}), each with distinct motion rules:
    foam is advected at the local fluid velocity and dissolves over time
    (\loc{FluidSim.compute}{547--553}); bubbles rise via a buoyancy term and
    relax toward local fluid velocity (\loc{FluidSim.compute}{555--561}); and
    spray falls under gravity with quadratic drag
    (\loc{FluidSim.compute}{563--568}).
    \item \textbf{Lifetime compaction.} Dead particles are stream-compacted
    out via an atomic counter and a ping-pong buffer
    (\texttt{WhiteParticlesCompacted}), with the surviving count and buffer
    copied back for the next frame (\loc{FluidSim.compute}{572--607}).
\end{itemize}

This is a complete, correctly-attributed (if uncredited in-file) 3D
implementation of a well-known SPH post-processing technique, but it is
visual/decorative: it reads the fluid simulation's state and does not feed
back into the pressure/density/viscosity solve (other than the base fluid
kernels it reuses for neighbor search).

\section{Volumetric Density Resampling}

\texttt{UpdateDensityTexture} (\loc{FluidSim.compute}{478--488}) resamples
\texttt{CalculateDensitiesAtPoint} on a regular voxel grid sized to the
container's aspect ratio (\loc{FluidSim.cs}{310--321}, driven by
\texttt{renderToTex3D}/\texttt{densityTextureRes},
\loc{FluidSim.cs}{41--42}) and writes the scalar density field into an
\texttt{RWTexture3D} for external volumetric rendering (e.g.\ ray marching or
isosurfacing, not present in this folder). This reuses the same summation-
density kernel already covered in Section 4, evaluated at arbitrary points
rather than at particle positions --- a direct application of the
continuous-reconstruction idea in study notes Section 4.2
(\emph{How Kernel Smoothing Produces a Continuous Approximation}, study
notes lines 551--775), $\langle A(\mathbf{x})\rangle \approx \sum_j m_j
W(\mathbf{x}-\mathbf{x}_j,h)$ (study notes line 767), evaluated off-particle.

\section{Summary Table}

\begin{longtable}{@{}p{3.6cm}p{5.6cm}p{5.8cm}@{}}
\toprule
\textbf{SPH Aspect} & \textbf{Implementation} & \textbf{Location} \\
\midrule\endhead
Neighbor search & Uniform grid + spatial hash + GPU counting sort & \file{SpatialHash3D.hlsl}:1--68; \file{FluidSim.compute}:146--185; \file{Helpers/SpatialHash/SpatialHash.cs}:30--36 \\
Density kernel & Spiky power-2 (non-standard choice, anti-clustering) & \file{FluidMaths3D.hlsl}:39--47, 69--73 \\
Near-density kernel & Spiky power-3 & \file{FluidMaths3D.hlsl}:28--36, 75--83 \\
Viscosity kernel & Poly6 (used for XSPH, not Laplacian viscosity) & \file{FluidMaths3D.hlsl}:17--26 \\
Density summation & $\rho_i=\sum_j W_{ij}$, unit mass assumed & \file{FluidSim.compute}:187--226 \\
Pressure closure & Linearized equation of state $p=(\rho-\rho_0)k$ & \file{FluidSim.compute}:85--88 \\
Near-pressure (stabilization) & Clavet-style second pressure field & \file{FluidSim.compute}:90--93, 267--346 \\
Pressure force & Symmetrized pairwise force & \file{FluidSim.compute}:326--332 \\
Viscosity & XSPH velocity smoothing & \file{FluidSim.compute}:409--459 \\
Time integration & Semi-implicit Euler + predicted positions + fixed sub-steps & \file{FluidSim.compute}:133--144, 461--475; \file{FluidSim.cs}:284--308 \\
Boundary handling & Kinematic box clamp + damped reflection & \file{FluidSim.compute}:101--131 \\
Foam/spray/bubbles & Ihmsen et al.\ 2012 classification \& transport & \file{FluidSim.compute}:356--388, 490--607 \\
Volumetric field export & Density resampled to 3D texture & \file{FluidSim.compute}:478--488 \\
\bottomrule
\end{longtable}

\section{Not Implemented / Possible Improvements}
\label{sec:notimplemented}

The following SPH mechanisms are discussed in the study notes (or are
standard extensions in the wider literature) but have no corresponding code
in \texttt{Simulation}:

\begin{itemize}
    \item \textbf{Explicit particle mass.} Study notes lines 758, 811, and
    822--827 all carry a per-particle mass $m_j$ through every kernel sum.
    The implementation has no mass field in any buffer or kernel; it is
    implicitly $1$ everywhere. This only happens to be correct because the
    spawner (\file{Spawner3D.cs}) places particles on a uniform grid of
    constant spacing, so variable-resolution, variable-density, or
    multi-material setups cannot be represented without first adding a
    stored mass (or volume) value and threading it through
    \texttt{CalculateDensitiesAtPoint} and \texttt{CalculatePressureForce}.
    \item \textbf{Full Tait equation of state / Mach-number-based sound
    speed.} Study notes lines 262--271 and 433--436 describe choosing $c_0$
    from a target Mach number and using exponent $\gamma$ (commonly 7 for
    water). The code's \texttt{pressureMultiplier} /
    \texttt{targetDensity} are free-standing stiffness constants with no
    physical sound-speed derivation, so density-error vs.\ stiffness
    trade-offs (study notes lines 433--436) cannot currently be reasoned
    about quantitatively; exposing an equivalent $c_0$ and Mach-number
    readout would make tuning more principled.
    \item \textbf{Incompressible/iterative pressure solvers (ISPH, PCISPH,
    DFSPH).} Study notes lines 271--292 and the Cummins--Rudman and
    Solenthaler--Pajarola references (study notes lines 947--949, 973--977)
    describe enforcing $\nabla\cdot\mathbf{v}=0$ via a pressure Poisson
    equation or predictive-corrective iteration. The simulation is purely
    weakly compressible with a single-pass explicit pressure force, so it is
    more prone to timestep-limited stiffness and audible/visible
    compressibility artifacts at high stiffness; adopting a PCISPH-style
    predict/correct loop (already partially hinted at by the existing
    ``predicted position'' mechanism in Section~\ref{sec:timeintegration})
    would be a natural incremental upgrade.
    \item \textbf{Physical viscosity discretization.} Study notes lines
    224--232 give the Navier--Stokes viscous term
    $\mu\nabla^2\mathbf{v}$, and the Morris--Fox--Zhu (1997) reference
    (study notes lines 940--944) gives a concrete SPH Laplacian for it. Only
    the non-physical XSPH smoothing (Section~\ref{sec:viscosity}) is
    implemented, so \texttt{viscosityStrength} does not correspond to a
    measurable kinematic viscosity $\nu$ and cannot be validated against a
    Reynolds-number-based benchmark (study notes lines 374--378).
    \item \textbf{Artificial viscosity ($\alpha,\beta$) and $\delta$-SPH
    density diffusion.} Study notes lines 452--463 list both as standard
    stabilizers for shocks and density/pressure noise. Neither is present;
    the only anti-noise/anti-clustering mechanisms are the near-pressure
    term (Section~\ref{sec:nearpressure}) and the ad hoc airborne-particle
    drag (\loc{FluidSim.compute}{351--355}).
    \item \textbf{Surface tension.} Study notes Section 3.1 (lines
    384--386) and the momentum equation (study notes line 232,
    $\mathbf{f}_s$) include an interfacial surface-tension force, and it is
    explicitly one of the features of M\"uller et al.\ (2003) cited in the
    reading list (study notes lines 950--955). No cohesion or curvature
    force of any kind is computed in \file{FluidSim.compute}; free-surface
    behavior is governed solely by the near-pressure anti-clustering term,
    so droplet break-up/cohesion behavior cannot be tuned independently of
    overall fluid stiffness.
    \item \textbf{Boundary/ghost particles and no-slip walls.} As discussed
    in Section~\ref{sec:boundary}, only a kinematic box clamp exists; there
    is no SPH-consistent wall treatment such as Adami--Hu--Adams (2012)
    (study notes lines 988--992), so wall friction, non-box geometry, and
    moving/rotating solid obstacles cannot currently be modeled without a
    substantial boundary-handling addition.
    \item \textbf{CFL-based adaptive timestep.} Study notes lines 469--489
    give an explicit formula for $\Delta t$ from $h$, $c_0$, $a_{\max}$, and
    $\nu$. The code only divides a user-chosen frame time by a fixed
    \texttt{iterationsPerFrame} and clamps the frame time via
    \texttt{maxTimestepFPS} (Section~\ref{sec:timeintegration}); it never
    measures per-particle velocity or acceleration to adapt the timestep,
    so a sudden high-energy event (e.g.\ a large splash) is not automatically
    guarded against beyond whatever margin the fixed sub-step count happens
    to provide.
    \item \textbf{Dead/unused interaction and edge-force uniforms.} As noted
    in Section~\ref{sec:boundary}, \texttt{edgeForce}, \texttt{edgeForceDst},
    \texttt{interactionInputPoint}, \texttt{interactionInputStrength}, and
    \texttt{interactionInputRadius} are declared in
    \file{FluidSim.compute}:44--53 but never read. If interactive
    mouse/controller forcing or a soft edge-repulsion force is desired (both
    common in SPH demos, including presumably a 2D counterpart to this
    project), the shader-side uniforms already exist and only need a kernel
    that consumes them plus CPU-side code to populate them per frame.
    \item \textbf{Adaptive/variable smoothing length.} Study notes lines
    408--416 discuss $h=\eta\Delta x$ and adaptive $h_i$ for variable
    resolution. \texttt{smoothingRadius} is a single global scalar
    (\loc{FluidSim.cs}{21}); there is no per-particle $h_i$ or refinement
    criterion.
    \item \textbf{Energy equation / thermal and multiphase effects.} Study
    notes Section 2.1.5 (lines 293--313) and the multiphase remarks in
    Section 2.1.6 are entirely out of scope here, which is reasonable for a
    single-phase isothermal graphics fluid but worth stating explicitly:
    there is exactly one fluid phase/material and no temperature field.
    \item \textbf{Convergence/validation tooling.} Study notes Section 3.7
    (lines 505--525) recommends refinement studies (halving $\Delta x$ and
    $\Delta t$, tightening solver tolerances) to confirm results are not
    parameter-dependent. No automated convergence test, benchmark scene, or
    regression harness exists in this folder; adding even a simple
    dam-break or hydrostatic-pressure test scene with a known analytic or
    reference result would make future parameter or kernel changes easier to
    validate.
\end{itemize}

\end{document}
